\documentclass[10pt]{article}
\usepackage[paperwidth=612pt, paperheight=792pt, left=180pt, right=54pt, top=36pt, bottom=45pt, headheight=14pt, headsep=18pt]{geometry}
\usepackage{amsmath,amssymb,amsfonts}
\usepackage{xcolor}
\usepackage{graphicx}
\usepackage{fancyhdr}
\usepackage[vietnamese]{babel}
\usepackage{newtxtext,newtxmath}

\definecolor{stewartcyan}{RGB}{0, 130, 185}

\pagestyle{fancy}
\fancyhf{}
\renewcommand{\headrulewidth}{0pt}
\fancyhead[L]{\textsf{\textbf{MỤC 2.5}}\hspace{1.5em}\textsf{Tính liên tục}}
\fancyhead[R]{\textsf{\textbf{123}}}

\setlength{\parindent}{14pt}
\setlength{\parskip}{0pt}

\begin{document}

\noindent
Một ứng dụng của Định lý giá trị trung gian là xác định vị trí nghiệm của phương trình như trong ví dụ sau đây.

\vspace{0.9em}
\noindent{\color{stewartcyan}\textsf{\textbf{VÍ DỤ 10}}}\hspace{1em}Chứng minh rằng phương trình
\[
4x^3 - 6x^2 + 3x - 2 = 0
\]
\noindent có một nghiệm nằm giữa 1 và 2.

\vspace{0.7em}
\noindent{\color{stewartcyan}\textsf{\textbf{LỜI GIẢI}}}\hspace{1em}Đặt $f(x) = 4x^3 - 6x^2 + 3x - 2$. Ta đang tìm một nghiệm của phương trình đã cho, nghĩa là một số $c$ nằm giữa 1 và 2 sao cho $f(c) = 0$. Do đó, ta lấy $a = 1$, $b = 2$, và $N = 0$ trong Định lý 10. Ta có
\begin{flalign*}
&& f(1) &= 4 - 6 + 3 - 2 = -1 < 0 &\\
\text{và} && f(2) &= 32 - 24 + 6 - 2 = 12 > 0 &
\end{flalign*}
Như vậy $f(1) < 0 < f(2)$; nghĩa là $N = 0$ là một số nằm giữa $f(1)$ và $f(2)$. Hàm số $f$ liên tục vì nó là một hàm đa thức, do đó theo Định lý giá trị trung gian tồn tại một số $c$ nằm giữa 1 và 2 sao cho $f(c) = 0$. Nói cách khác, phương trình $4x^3 - 6x^2 + 3x - 2 = 0$ có ít nhất một nghiệm $c$ trong khoảng $(1, 2)$.

\vspace{0.5em}
Thực tế, chúng ta có thể định vị nghiệm một cách chính xác hơn bằng cách áp dụng lại Định lý giá trị trung gian. Vì
\[
f(1{,}2) = -0{,}128 < 0 \qquad\text{và}\qquad f(1{,}3) = 0{,}548 > 0
\]
\noindent nên nghiệm phải nằm giữa $1{,}2$ và $1{,}3$. Bằng cách thử và sai, máy tính cầm tay cho kết quả
\[
f(1{,}22) = -0{,}007008 < 0 \qquad\text{và}\qquad f(1{,}23) = 0{,}056068 > 0
\]
\noindent do đó nghiệm nằm trong khoảng $(1{,}22;\, 1{,}23)$.\quad{\color{stewartcyan}\rule{1.15ex}{1.15ex}}

\vspace{0.8em}
Chúng ta có thể sử dụng máy tính bỏ túi có chức năng vẽ đồ thị hoặc máy tính để minh họa việc áp dụng Định lý giá trị trung gian trong Ví dụ 10. Hình 10 cho thấy đồ thị của $f$ trong khung nhìn $[-1, 3] \times [-3, 3]$ và bạn có thể thấy rằng đồ thị cắt trục hoành tại một điểm giữa 1 và 2. Hình 11 cho thấy kết quả khi phóng to vào khung nhìn $[1{,}2;\, 1{,}3] \times [-0{,}2;\, 0{,}2]$.

\vspace{0.8em}
\begin{center}
\begin{minipage}[t]{0.47\textwidth}
    \centering
    \includegraphics[width=\linewidth]{images/p0158_fig1.png}\\[5pt]
    {\small\textsf{\textbf{HÌNH 10}}}
\end{minipage}\hfill
\begin{minipage}[t]{0.47\textwidth}
    \centering
    \includegraphics[width=\linewidth]{images/p0158_fig2.png}\\[5pt]
    {\small\textsf{\textbf{HÌNH 11}}}
\end{minipage}
\end{center}

\vspace{0.8em}
Thực ra, Định lý giá trị trung gian đóng một vai trò trong chính cách thức hoạt động của các thiết bị vẽ đồ thị này. Máy tính sẽ tính toán một số hữu hạn các điểm trên đồ thị và làm sáng các điểm ảnh (pixel) chứa những điểm đã tính toán đó. Nó giả định rằng hàm số liên tục và nhận tất cả các giá trị trung gian giữa hai điểm liên tiếp. Do đó, máy tính ``nối các điểm'' lại bằng cách làm sáng các điểm ảnh ở giữa.

\end{document}